\section{硬件、数据与展示产物}
\subsection{硬件与控制接口}
\begin{table}[htbp]
\centering\small
\caption{现有基础设施及待补配置。}
\begin{tabularx}{\linewidth}{@{}lY@{}}
\toprule
项目 & 当前记录\\\midrule
机器人 & 单臂 UR5e\\
相机 & 固定主视角、腕部视角；型号\pending\\
深度 & 拟保存主视角 RGB 对齐深度\\
机器人状态 & 关节位置／速度、TCP 位姿\\
夹爪 & 指令状态；具体硬件\pending\\
采集方式 & Freedrive 示范；实际采样率\pending\\
动作接口 & 单臂映射与执行频率\pending\\
推理硬件 & GPU、延迟与动作块长度\pending\\
\bottomrule
\end{tabularx}
\end{table}

\begin{figure}[htbp]
\centering
\panel{31mm}{硬件照片占位\\标注 UR5e、夹爪、固定相机、腕部相机\\显示桌面工作区域与相机朝向}
\caption{硬件与视角布置。参考 FlowWAM Figure 6 和 OpenWAM Figure 18 的标注方式。}
\end{figure}

\subsection{原始记录与离线处理}
已有采集代码保存双视角 RGB、RTDE 状态、夹爪指令状态、episode 元数据和同步索引。新增记录项包括主视角深度、深度单位、RGB--depth 对齐方式、相机内参、相机帧号与时间戳。

离线处理生成前景 mask、机器人／物体角色标签，以及结合 RGB 光流和深度得到的 Track UVD。处理版本、有效像素比例和异常帧随 episode 保存。固定主视角作为首选 Track 视角；腕部视角提供 RGB 观察。

\subsection{相机参数与标定路径}
RealSense SDK 提供出厂内参、同一相机内部的彩色／深度外参和深度尺度；这些参数可直接读取保存。相机到机器人基座的变换属于另一层参数，涉及世界坐标点云融合、机器人投影或移动相机运动补偿时再配置。当前以固定相机中的 UVD 表示组织几何数据。

\subsection{训练、部署和展示}
分别记录训练监督、部署输入和生成输出：训练可以使用未来深度与轨迹；部署只使用当前可获得的观察；预测分支输出作为内部表示和可选可视化。记录视频、Track、动作专家的训练范围，以及推理时各分支的实际采样配置。

实机执行录像、模型生成视频和轨迹预测分别保存并标注。展示策略执行时使用实机录像；展示预测质量时使用对应的预测和参考片段。

\subsection{交付材料清单}
主文保留四任务场景图和一张含 OOD 的主结果表。附录包括硬件照片、任务终点表、随机化范围、示范／评估协议和数据处理说明。可选材料包括 OOD 对照图、预测与执行序列、失败案例和推理开销。

\paragraph{基础设施版本。}\href{https://github.com/Yutenji-Nyamu/UR5e-RoboTwin-Real}{UR5e-RoboTwin-Real}，本次只读核对的提交为 \texttt{097922c}。该版本的相机采集流为 RGB，深度记录列为本次方案的待实现项。
